\section{Du registre de trajectoire à l’opérateur de masse
}
\label{vi:sec:operador-masa}

La construction distingue deux opérations. D’abord, la dynamique produit
un registre de la trajectoire. Ensuite, une fonctionnelle extrait des coordonnées
entières et les évalue au moyen des constantes internes déjà obtenues.
Conserver cet ordre explique comment la masse peut être une grandeur de la
persistance sans épuiser sa structure.


\subsection{Lecture chronologique et événements orientés
}

Soit \(d_1,\ldots,d_{108}\) la lecture numérique d’une route enrichie.
Dans la coordinatisation déclarée, le support d’action est \(\{1,3,5,7\}\), et

\[
 n_A=\#\{t:d_t\in\{1,3,5,7\}\}.
\]
On conserve en outre les \(12\) événements des blocs de \(9\) pas :

\[
 e_m=\sigma_m\#\{t:9m+1\le t\le9m+9,\ d_t\in\{2,4,6,8\}\},
 \quad 0\le m<12,
\]
où \(\sigma=(+,+,+,-,-,-,+,+,+,-,-,-)\). Le support de cet
extracteur est une partie explicite de la coordinatisation de lecture ; sa fonction ne
doit pas être confondue avec le sélecteur ternaire de phase du TPK.


La distribution des événements est pertinente même lorsque leurs sommes
coïncident. Pour la voir, nous introduisons

\[
 p_i=e_i+e_{i+6},\quad a_i=e_i-e_{i+6},\quad
 s_C=\sum_{i=0}^{5}p_i,\quad M_i=p_i-p_5\ (0\le i<5).
\]

\begin{proposition}[Reconstruction entière des événements
]
\label{vi:prop:eventos}
La transformation \((e_0,\ldots,e_{11})\mapsto(s_C,M_0,\ldots,M_4,a_0,\ldots,a_5)\)
est injective sur \(\mathbb Z^{12}\). Son image est caractérisée par

\[
 s_C-\sum_{i=0}^4M_i\equiv0\pmod6,
 \qquad p_i\equiv a_i\pmod2\quad(0\le i<6),
\]
où \(p_5=(s_C-\sum M_i)/6\) et \(p_i=p_5+M_i\) pour \(i<5\).

\end{proposition}
\begin{proof}
Sommer les six expressions de \(p_i\) donne
\(s_C=6p_5+\sum_{i<5}M_i\). Cela récupère les \(p_i\), puis

\[
 e_i=\frac{p_i+a_i}{2},\qquad
 e_{i+6}=\frac{p_i-a_i}{2}.
\]
Les congruences sont nécessaires pour que ces quantités soient entières et
suffisantes pour que les formules définissent une préimage entière. Substituer
dans la transformation directe récupère toutes les coordonnées.

\end{proof}

La décomposition \(12=1+5+6\) acquiert ainsi un sens opératoire :
le total, les différences et les orientations récupèrent les événements.
La division par \(6\) appartient à cette reconstruction. Lorsque le caractère
utilise \(s_C/6\), \(s_C\) reste la somme entière ; une coordonnée déjà
divisée n’est pas divisée une seconde fois.


\subsection{Mémoire torsionnelle et signature
}

Pour chaque occurrence de \(9\), la lecture conserve son prédécesseur.
Nous définissons, avec des indices cycliques,

\[
 j_2(d)=\begin{cases}0,&d=9,\\1,&d\in\{2,4,6,8\},\\-1,&d\in\{1,3,5,7\},\end{cases}
 \quad z_t=\mathbf1_{\{d_t=9\}}j_2(d_{t-1}).
\]
Alors
\[
 T_z=\sum_{t=1}^{108}z_t,\quad
 C_z=\sum_{t\in\{27,54,81,108\}}z_t,\quad
 k_{\Delta_4}=T_z-2C_z.
\]
La première somme conserve la mémoire accumulée et la seconde distingue les
fermetures de couronne. Éliminer les prédécesseurs avant d’évaluer ces
expressions altère l’objet lu.


Avec \(\tau_m=\operatorname{sgn}(e_m)\), les autres coordonnées sont

\[
 b_{120}=\sum_{a=0}^{3}\operatorname{sgn}(e_a+e_{a+4}+e_{a+8}),
 \qquad b_{\zeta}=\sum_{m=0}^{11}\tau_m\tau_{m+3},
\]
où le second indice est modulo \(12\). On obtient la signature

\[
 L_{\mathrm{tick}}(\gamma)
 =\sigma(\gamma)=(n_A,s_C,k_{\Delta_4},b_{120},b_{\zeta})\in\mathbb Z^5.
\]
Cette application est déterministe dans la coordinatisation déclarée. Les signes et les
corrélations expliquent pourquoi la composition doit se réaliser sur le
registre avant d’extraire les \(5\) entiers : la somme des signatures ne
représente pas, en général, la composition de deux histoires.


\subsection{Caractère et spécialisation interne
}

Soit \(X=(X_A,X_C,X_4,X_{120},X_{270})\). Le caractère formel est

\[
 \chi(\sigma;X)=\exp\!\left(n_AX_A+s_CX_C+k_{\Delta_4}X_4+
 b_{120}X_{120}+b_{\zeta}X_{270}\right).
\]
Pour chaque \(X\), il satisfait
\(\chi(\sigma+\sigma';X)=\chi(\sigma;X)\chi(\sigma';X)\)
parce que l’exposant est linéaire en la signature. C’est une loi sur les signatures ; elle
n’affirme pas que l’extracteur est additif sur les routes sans conserver l’interface.


La spécialisation emploie les coordonnées internes construites précédemment :

\[
 X_A=A_{\mathrm{rad}},\quad X_C=C^*_{\mathrm{rad}}/6,
 \quad X_4=\Delta_4,\quad X_{120}=s_{120},\quad
 X_{270}=\zeta_{270}:=135\Delta_4^2.
\]
Le caractère est adimensionnel et positif. Son effet est une homothétie sur
la droite de masse. La forme exponentielle détermine comment agissent les
coefficients ; leur provenance est fixée dans les registres et les constructions
antérieurs.


\subsection{Hiérarchie bicouche et conservation des deux généalogies
}

Nous écrivons \(x=A_{\mathrm{rad}}\), \(y=C^*_{\mathrm{rad}}\). Dans la
chambre \(x>y>0\), les canaux sont

\[
 q_+=e^{-x-y},\qquad q_-=e^{-x+y},\qquad0<q_+<q_-<1.
\]
Pour une involution \(R\), avec \(P_\pm=(I\pm R)/2\), l’opérateur
des deux canaux est

\[
 T=q_+P_++q_-P_-.
\]
Ses moments scalaire et orienté s’expriment sous la forme

\[
 c_n=q_-^n+q_+^n,\qquad s_n=q_-^n-q_+^n.
\]
Multiplier ces expressions prouve

\[
 c_n^2-s_n^2=4e^{-2nx},\quad
 c_{m+n}=\frac{c_mc_n+s_ms_n}{2},\quad
 s_{m+n}=\frac{s_mc_n+c_ms_n}{2}.
\]
Les deux suites satisfont

\[
 X_{n+2}=c_1X_{n+1}-e^{-2x}X_n,\qquad X\in\{c,s\}.
\]
Ces relations montrent que les hauteurs partagent un même opérateur.
On ne définit pas de collection différente de moments pour chaque espèce.


\begin{proposition}[Semi-groupe des hauteurs
]
\[
 \langle90,120\rangle=\{90,120\}\cup\{30r:r\ge6\}.
\]
\end{proposition}
\begin{proof}
Après division par \(30\), il suffit d’étudier \(\langle3,4\rangle\).
Les entiers \(6,7,8\) sont \(3+3,3+4,4+4\) ; ajouter \(3\) produit
tous les entiers suivants. En dessous de \(6\), les seuls
éléments positifs sont \(3,4\).

\end{proof}

Les hauteurs \(360=4\cdot90=3\cdot120\) illustrent la différence entre
égalité d’évaluation et mémoire de composition. Le quotient algébrique
identifie les deux exposants ; le registre conserve si le transport
a employé quatre étapes d’un type ou trois de l’autre.


La spécialisation raffinée a pour exposant

\[
 \log\chi_{\mathrm{ref}}=n_Ax+\frac{s_C}{6}y+k_{\Delta_4}\Delta_4
 +b_{\zeta}\zeta_{270}+\sum_{h\in\langle90,120\rangle}N_hs_h,
 \qquad N_h=\eta_h+b_{120}\mathbf1_{\{h=120\}}.
\]
La paire \((\eta_{120},b_{120})\) est conservée même si l’évaluation utilise
sa somme. On conserve également séparés \(s_{270}\) et
\(\zeta_{270}\) : le premier est un moment bicouche et le second une
expression quadratique du raffinement.


\begin{proposition}[Convergence de l’évaluation raffinée
]
Si \(\sum_h|N_hs_h|<\infty\), le caractère raffiné définit un scalaire
positif fini et ses évaluations tronquées convergent vers lui.

\end{proposition}
\begin{proof}
L’hypothèse donne la convergence absolue de l’exposant. La continuité de
l’exponentielle permet de passer à la limite ; l’exponentielle d’un réel fini est
strictement positive.

\end{proof}

La convergence évalue une suite de coefficients déjà spécifiée.
La sélection de ces coefficients par la route est une opération différente,
et doit accompagner toute évaluation attribuée à une espèce.


La notation \(\chi_{\mathrm{full}}\) conservée dans certaines formules
d’évaluation des secteurs désigne ce même caractère
\(\chi_{\mathrm{ref}}\), avec les mêmes arguments et coefficients.
Le changement de symbole n’ajoute ni autre spécialisation ni autre tour.


\subsection{Lecteurs de déplacement et d’autocorrélation
}

La même route admet deux lectures du retour :

\[
 K_D=\left(D_{27}+D_{36},D_1,\frac{D_4-D_3}{2}\right),\qquad
 K_\Omega=(\Omega,54,P_6).
\]
La première application conserve les défauts de déplacement ; la seconde, la mémoire de
retour et l’autocorrélation. Dans le recensement complet de \(324\) orientations,
il y a \(14\) profils de la première, \(9\) de la seconde et \(40\) paires
conjointes. Les \(18\) coïncidences ont pour valeur \((80,54,6)\).
On n’en infère pas l’égalité des deux opérateurs sur tout l’atlas.


Pour la route électronique, les déplacements sont

\[
 (D_1,D_3,D_4,D_{27},D_{36})=(54,56,68,48,32).
\]
La substitution produit \(K_D=(80,54,6)\), coïncidant avec
\(K_\Omega\). Son exposant est

\[
 \kappa_e=80-54\alpha+6\Delta_4.
\]
La trajectoire électronique est construite dans
\eqref{vi:dep:trayectoria-electron} ; ses lecteurs de déplacement et
d’autocorrélation sont évalués, respectivement, dans
\eqref{vi:dep:KD} et \eqref{vi:dep:KOmega}. La coïncidence
\(K_D=K_\Omega=(80,54,6)\) appartient au théorème de fermeture bêta de la
sous-section~\ref{vi:dep:beta}. Ici, cette exception centrale est conservée et
l’opération s’étend aux fibres de l’atlas.


\subsection{Opérateur sur une fibre et section dimensionnelle
}

Soit \(\mathscr H_M\) l’espace de Hilbert fini de base orthonormale
\((e_\gamma)\) étiquetée par les routes d’une multisection. Pour
\(r\in\{D,\Omega\}\), l’opérateur de retour est défini par
\(B_M^re_\gamma=\kappa_\gamma^re_\gamma\), avec des exposants réels ;
il est donc autoadjoint. Soit \(M_M^{(0)}\) l’opérateur basal positif,
dont la connexion avec le caractère est fixée ci-dessous.
Le rapport adimensionnel des
sections d’action est \(R_{\mathrm{act}}>0\). On définit

\[
 M_M^{\beta,r}=R_{\mathrm{act}}^{B_M^r/2}
 M_M^{(0)}R_{\mathrm{act}}^{B_M^r/2}.
\]

\begin{proposition}[Positivité et restriction scalaire
]
L’opérateur \(M_M^{\beta,r}\) est positif. Si les deux opérateurs sont
diagonaux dans la base de routes, chaque valeur basale \(m_\gamma^{(0)}\)
se transforme en

\[
 m_\gamma^\beta=m_\gamma^{(0)}R_{\mathrm{act}}^{\kappa_\gamma}.
\]
\end{proposition}
\begin{proof}
Soit \(S=R_{\mathrm{act}}^{B_M^r/2}\), autoadjoint et inversible.
Pour tout \(v\),
\(\langle v,SM_M^{(0)}Sv\rangle=\langle Sv,M_M^{(0)}Sv\rangle\ge0\).
Dans la base diagonale, les deux facteurs \(S\) multiplient chaque valeur
par \(R_{\mathrm{act}}^{\kappa_\gamma/2}\) ; leur produit donne la formule.

\end{proof}

La section absolue de masse s’obtient à partir de l’action et de l’horloge :

\[
 \omega_0=\frac{2\pi}{108t_0},\qquad
m_{\mathrm{clk}}=\hbar_{\mathrm{ret}}\omega_0c_{\mathrm{int}}^{-2}.
\]
Dans la coordinatisation normalisée par rapport à l’électron, l’opérateur basal précédent est

\[
 M_M^{(0)}e_\gamma
 =m_e^{(0)}\,
   \chi_{\mathrm{ref}}(\sigma_\gamma-\sigma_e,
                       \eta_\gamma-\eta_e)e_\gamma.
\]
Ici, \(m_e^{(0)}\) est la réalisation de la coordonnée électronique
basale, avec son préfacteur, son raffinement et sa base dimensionnelle.
Si \(E_e^{(0)}=\varepsilon_e^{(0)}\mathcal U_E\), alors

\[
 m_e^{(0)}=\varepsilon_e^{(0)}\mathcal U_Ec_{\mathrm{int}}^{-2}
 =b_e\,m_{\mathrm{clk}},\qquad
 b_e=\varepsilon_e^{(0)}\frac{\mathcal U_E}{E_{\mathrm{clk}}}.
\]
Dans la coordinatisation \(\mathcal U_E=E_{\mathrm{clk}}\), on obtient
\(b_e=\varepsilon_e^{(0)}\) ; dans une autre coordinatisation, le quotient
de bases indiqué est conservé. Si les coordonnées de tour sont déjà
relatives à l’électron, on omet une seconde soustraction de \(\eta_e\).
Cette égalité relie l’opérateur à la signature de chaque route et conserve
explicitement la normalisation électronique. Un choix différent de
section de référence doit également transporter cette normalisation.
L’échelle de la constante réduite de Planck appartient à cette construction
dimensionnelle. Le caractère et le facteur de retour agissent ensuite comme
scalaires. L’électron conserve sa structure : il ne s’identifie pas au
quantum chronologique \(m_{\mathrm{clk}}\).


Pour deux routes dans la même coordinatisation absolue,

\[
 \frac{m_Y^\beta}{m_X^\beta}
 =\chi_{\mathrm{ref}}(\sigma_Y-\sigma_X,\eta_Y-\eta_X)
 R_{\mathrm{act}}^{\kappa_Y-\kappa_X}.
\]
La section dimensionnelle commune s’annule. La différence des exposants de
retour ne s’annule pas sauf si elle est nulle. Cette identité conserve la
distinction entre échelle absolue et structure relative d’une collection.


\subsection{Le secteur de propagation nulle
}

L’exponentielle prend des valeurs strictement positives. La coordinatisation nulle se
construit avant ce caractère, dans la réalisation déployée

\[
 \mathcal A_{\mathrm{sp}}=\mathbb R[R]/(R^2-1),\quad
 z=r_+P_++r_-P_-,\quad N(z)=r_+r_-.
\]
Les directions \(\mathbb RP_+\) et \(\mathbb RP_-\) ont une norme
nulle. Une propagation transportée dans l’une d’elles appartient à ce secteur ;
la lecture matérielle bicouche emploie le couplage des deux. De cette manière,
le photon et le secteur gluonique conservent leur place dans la théorie même
s’ils ne constituent pas une ligne du caractère positif. La réalisation de leurs
représentations est exposée avec les secteurs correspondants.

